
        \documentclass[fleqn]{article}      
        \usepackage{amsmath}
        \usepackage{fontspec}
        \usepackage{kotex} % 한국어 지원  

        \begin{document}      
        쉬움: \\\\  
1. 다음 중 정적분의 정의는 무엇인가요?  
   a) 특정 구간의 기울기를 구하는 것  
   b) 함수의 면적을 구하는 것  
   c) 함수의 최대값을 구하는 것  
   d) 함수의 최소값을 구하는 것 \\\\  

2. ∫2x dx의 적분 결과는 무엇인가요?  
   a) x^2 + C  
   b) 2x^2 + C  
   c) x^2  
   d) 2x + C \\\\  

3. 정적분 ∫(0부터 1까지) x dx의 값은 무엇인가요?  
   a) 0.5  
   b) 1  
   c) 2  
   d) 3 \\\\  

4. 다음 중 부정적분의 결과로 올바른 것은? ∫3x^2 dx  
   a) x^3 + C  
   b) 3x^3 + C  
   c) x^2 + C  
   d) x^3 \\\\  

5. 함수 f(x) = x^2의 그래프 아래 면적을 구하기 위해 사용하는 방법은?  
   a) 미분  
   b) 적분  
   c) 대수  
   d) 통계 \\\\  

중간: \\\\  
1. 함수 f(x) = 3x^2의 정적분 ∫(1부터 2까지) f(x) dx의 결과는?  
   a) 8  
   b) 9  
   c) 10  
   d) 11 \\\\  

2. ∫(2x + 3) dx의 부정적분 결과는 무엇인가요?  
   a) x^2 + 3x + C  
   b) x^2 + 3x  
   c) 2x^2 + 3x + C  
   d) 2x^2 + 3x \\\\  

3. 다음 중 정적분의 기하적 의미는 무엇인가요?  
   a) 함수의 기울기  
   b) 함수의 면적  
   c) 함수의 교차점  
   d) 함수의 대칭성 \\\\  

4. 함수 f(x) = x^3의 정적분 ∫(0부터 1까지) f(x) dx의 값을 구하세요. \\\\  

5. 주어진 함수 f(x) = 4x의 적분을 구하고, 그 결과를 설명하세요. \\\\  

어려움: \\\\  
1. ∫(x^2 + 2x + 1) dx의 부정적분 결과는 무엇인가요? \\\\  

2. 다음 정적분의 값을 구하세요: ∫(1부터 3까지) (2x - 1) dx. \\\\  

3. f(x) = 2x^3 - 4x의 정적분을 구하고, 그 결과를 설명하세요. \\\\  

4. 함수 f(x) = sin(x)의 정적분 ∫(0부터 π/2까지) f(x) dx의 값을 구하세요. \\\\  

5. 정적분 ∫(0부터 1까지) (x^4 - 2x^3 + x^2) dx의 결과는 무엇인가요? \\\\  

      
        \end{document} 
    